\begin{theorem}[A strict joint odd CM gate]
\label{thm:joint-odd-cm-gate}
Let \(T_{\rm odd}\) be the maximal pro-odd projection of the
corrected Eisenstein CM homomorphism, and put
\[
                       V_D=U_D\cap\ker T_{\rm odd}.
\]
Then \(V_D\) is closed and \(G\)-normal, and
\begin{equation}
 B_{\max2}\cong
 \left(V_D/\overline{[G,V_D](V_D\cap H)}\right)_{\max2}.
                                                               \label{eq:joint-odd-cm-1}
\end{equation}
The original restriction kernel is
\begin{equation}
 \ker\bigl(X_2\to H^1(V_D,\mathbf F_2)\bigr)
   =\bigoplus_{q\ {\rm odd}}\mathbf F_2\psi_q
                            \oplus\mathbf F_2\eta_3.    \label{eq:joint-odd-cm-2}
\end{equation}
Moreover
\begin{equation}
 [U_D:V_D]=\infty,\qquad
 [U_D\cap H:V_D\cap H]=\infty.                          \label{eq:joint-odd-cm-3}
\end{equation}
\end{theorem}

\begin{proof}
Semilinear conjugation acts on the cubic norm-one torus by
identity or inversion, and preserves its pro-odd factor.
Thus \(V_D\) is closed and \(G\)-normal, and
\(U_D/V_D=T_{\rm odd}(U_D)\) is pro-odd.
A \(G\)-invariant character of \(V_D\) with values in a finite
abelian two-group extends uniquely to \(U_D\), by the central
coprime averaging argument. If it kills \(V_D\cap H\), its
extension kills \(U_D\cap H\), since the quotient of these
intersections is pro-odd.
Hence the relative maximal pro-two abelian quotients agree,
giving \eqref{eq:joint-odd-cm-1} from the established reduction on \(U_D\).
A character of \(U_D\) with values in \(\mathbf F_2\) and
vanishing on \(V_D\) also vanishes, so the original
restriction kernels agree and give \eqref{eq:joint-odd-cm-2}.
No factorization \(G=U_DH\) is used.

For strictness use the arithmetic image of the corrected
Tate representation of
\[
                         E_3:y^2=x^3+1.
\]
The established cubic CM image is \(N_3(4)\).
At the split prime thirteen,
\[
 \mathbf Z[\zeta_3]\otimes\mathbf Z_{13}
           \cong\mathbf Z_{13}\times\mathbf Z_{13},
\]
choose the norm-one element \(v=(14,14^{-1})\), of infinite
pro-thirteen order, and put the identity at every other
component. The level-four condition at two is satisfied.
Choose \(m\in M\cap H\) with this corrected cubic CM value,
and put \(a=m^{\varepsilon_{13}}\).
Then \(a\) is pro-thirteen and has full multiplier one.
Its cubic Tate value is \(v\): the central
two-primary correction in the corrected representation is
killed by the power \(\varepsilon_{13}\).

Choose \(g\in H\) with full multiplier
\[
 \lambda_3(g)=-1,\qquad
 \lambda_p(g)=1\quad(p\ne3),
\]
and let \(h=g^{\varepsilon_2}\).
It is pro-two with the same multiplier tuple.
For every odd prime \(p\ne3\), its action \(\Phi_p(h)\)
is trivial, since it is both pro-two and contained in the
pro-\(p\) image of \(\ker\lambda_p\).
For \(p\ne13\), the normalized action \(\Phi_p(a)\) is
trivial because \(a\) is pro-thirteen.
Thus the commutator
\[
                              w=[h,a]
\]
has trivial \(\Phi_p\)-image for every prime \(p\):
at thirteen use the first entry, and at all other primes
use the second.
Its rational Kummer coordinates vanish as well, because
those of \(a\) are supported at thirteen, where
\(\lambda_{13}(h)=1\):
\[
 \rho_q^I(w)
       =(\lambda(h)-1)\rho_q^I(a)=0.
\]
Consequently \(w\in P_D\cap H\).
The cubic defect vanishes arithmetically on \(D\), so
\(w\in T_D\cap H\).

The element \(h\) fixes \(\mathbf Q(i)\), because its
dyadic multiplier is one. Its harmonic Tate action
is therefore linear over the CM coefficient ring and
commutes with that of \(a\). The harmonic value of \(w\),
and in particular \(W_{\rm odd}(w)\), is one.
Hence \(w\in U_D\cap H\).

On the other hand, \(h\) conjugates
\(\mathbf Q(\sqrt{-3})\), since its multiplier at three
is minus one. It inverts the cubic norm-one torus.
The Tate commutator has value \(v^{-2}\);
the central two-primary correction disappears under the
odd projection. Thus
\[
                         T_{\rm odd}(w)=v^{-2},
\]
which has infinite pro-thirteen order.
The procyclic closure of \(w\) therefore has infinite
image in \(U_D/V_D\), proving both assertions in \eqref{eq:joint-odd-cm-3}.
\end{proof}

\begin{remark}
This witness deliberately uses a prime distinct from two
and three. It proves a strict joint gate but makes no claim
that the entire odd cubic CM image survives on \(U_D\).
In particular the wild three-primary part requires a
separate analysis. The theorem does not assert \(G=V_DH\).
\end{remark}
